\paragraph{Monochromatic cycles in the posting graph.}
Draw an edge $j\xrightarrow{r}u$ whenever a negative row at $j$ posts
$u$ at $r$. A cycle is monochromatic if all its edges have the same
color. Individual rows may post several colors.

\begin{lemma}[Backward induction on the input names]
\label{core:lem:ranked-posting-survivor}
Choose one row $t^j$ at each input $j$. If their posting graph has no
monochromatic cycle, every word supplies at least one head's mark on
its chosen row. Conversely, the fold on any monochromatic cycle, taken
at its color, omits every chosen mark. Whenever there is an omitting
word, its minimum possible number of leaves equals the length of a
shortest monochromatic cycle.
\end{lemma}
\begin{proof}
Give each owner $r$ a rank on the inputs, strictly decreasing along
the $r$-edges. Evaluate a word to an input name by choosing at each
$r$-gate the child whose resulting input has smaller $r$-rank,
breaking ties arbitrarily. If a subtree wins the $r$-test on a set
$E$, its resulting input has $r$-rank at most the largest rank in
$E$. This follows by induction: an $r$-gate can choose at least as
well as a winning child, and at every other gate both children win.

Let $j$ be the resulting input at the root. Follow the selected path
to its leaf; every selected child on this path also has result $j$.
An $r$-sibling cannot win on $t^j_r$: all its postings have rank
strictly below $j$, whereas the gate chose $j$. On a probe with full
menu at $j$ no sibling on the path can therefore absorb at its gate.
The leaf supplies the mark. This argument also works on an infinite
input set: restrict the rows and the ranks to the finitely many
names in the word. Conversely, in a cycle fold every used
head has a cycle successor posted at the fold owner and every other
head is absent. Restricting the obstruction to the support of an
omitting word gives a cycle within its leaves; a shortest cycle uses
each of its vertices once.
\end{proof}

For comparison, if all losses at a head use a single designated owner
$c(j)$, there is a direct extraction even with several losses there.
Assign any owner to inputs without losses.

\begin{lemma}[A single-owner fold in every feasible support]
\label{core:lem:monochromatic-safe-fold}
The support of every word omitting these losses contains a nonempty
set $S$ of inputs of one designated owner $r$ whose $\star_r$-fold
also omits them. In particular every inclusion-minimal feasible
support has a read-once realizer of this form.
\end{lemma}
\begin{proof}
Write $H_r=\{j:c(j)=r\}$. Some $r$ has
$\Phi_r(V\setminus H_r)=0$, since otherwise the positive coalitions
$V\setminus H_r$ would have empty intersection, contrary to
\cref{core:thm:profiles}. Keep both children of each $r$-gate and
one child losing this test at every foreign gate. The resulting
$r$-fold has all its leaves in $H_r$. Test a pure $r$-row with the
full menu at its head: the source bit is independent of the other
winners in that menu by source decomposition. Every foreign gate
then acts by union, and $\star_r$ is monotone, so this pruning
preserves omission there. On a foreign row its head is
absent from the fold. Idempotence removes repeated inputs.
\end{proof}

\begin{theorem}[A cut with one loss per head]
\label{core:thm:monochromatic-head-cut}
Suppose every head with several losses belongs to monochromatic cycles
of at most one color. Let $(i,p)$ be a positive query for which no
vertex of an $r$-cycle is posted at $r$ in $p$. These hypotheses
allow arbitrary mixed rows. The first is automatic when each head
has one loss, or all its losses use only its designated owner.
The query is impossible exactly when one
can choose a set $H\ni i$ and one loss $(j,t^j)$ at each $j\in H$
such that
\[
 t^j_r\setminus H\subseteq p_r\quad(j\in H,\ r\in I),\qquad
 \text{the chosen rows on $H$ have no monochromatic cycle}.
\]
The query is decidable in polynomial time from either the tails or
the canonical prime guards. A possible query has a witness with at
most $1+K|V|^2$ leaves.
\end{theorem}
\begin{proof}
For the obstruction direction restrict to the chosen rows.
By \cref{core:lem:unmarked-background}, a witness would give a word
on $H$ omitting every restricted loss, contrary to
\cref{core:lem:ranked-posting-survivor}. This direction requires
neither hypothesis.

For the converse, on the remaining rows take the greatest support
$S_r\subseteq V\setminus p_r$ of a safe $r$-fold, allowing empty
support. Start with $V\setminus p_r$ and remove an input whenever
one of its rows has no $r$-posting in the current set. A fold must
absorb whenever it uses the marked head, so this computes $S_r$.
Delete the rows with an $r$-posting in $S_r$. The fold omits all
remaining marks and cannot absorb at the positive. Attaching it by
$\star_r$ therefore lifts any witness after deletion to one before
deletion. Repeat over the owners. If $i$ no longer heads a row,
its projection succeeds.

At stabilization let $H$ be the remaining heads, with $i\in H$.
Every tail posting outside $H$ agrees with $p$, since its unheaded
projection could otherwise have deleted the row. For each $r$ start
with the heads on $r$-cycles in the remaining graph restricted to $H$.
Iteratively remove any head having a remaining row whose $r$-postings
miss the current set. If a nonempty set remains, choose one $r$-posting
in that set from every row at each surviving head. A sink strongly
connected component has an edge out of each vertex. Every vertex
therefore lies on a monochromatic cycle and is outside $p_r$.
The component meets every tail at its heads, so its fold is safe
and hits a remaining row. This contradicts stabilization. Thus every
head in the starting set is removed. At each such head choose its
removing row. The choices for different colors agree wherever they
overlap, since a head lying on cycles of two colors has just one row.
Choose any remaining row at every other head. An $r$-cycle would
consist entirely of heads from the $r$-search, and each edge on it
leads to an earlier removal. This proves the required acyclicity.

To implement deletion, record $U_r$, the union of the fold supports
used at $r$. A row remains when its $s$-postings miss $U_s$ for
every $s$. In the support test, further exclude the $r$-postings
on the current candidate $S_r$. A head can remain in $S_r$ exactly
when no zero survives these exclusions. Scan the listed rows, or
on each observer facet test whether the allowed box is positive by
containment of a listed prime in its set of exclusions. Each strict
increase of a register adds an input-owner pair. At most $K|V|$
folds, each with at most $|V|$ leaves, and a final projection suffice.
\end{proof}

The same criterion handles prescribed singleton tests when the
augmented graph satisfies its hypotheses. For a required win $s:A$,
add the pure $s$-row posting $A$ at each head outside $A$, as in
\eqref{core:eq:total-monochromatic-source}. For a required loss
$\Phi_r(B)=0$, some input outside $B$ must supply a mark when just
$B$ is posted at $r$: give every unposted input the full menu and
apply source decomposition. Thus one may test each required mark
and each losing test separately on the augmented rows, and take
the union of their witnesses. This preserves every required singleton
win and every omission. If there are only winning tests, any feasible
word has a source at the empty coloring.

With just one loss at each headed input this is a reachability test.
From $i$ follow the edges $j\xrightarrow{r}u$ with $u\notin p_r$.
The mark is possible exactly when this search reaches an unheaded
input or a monochromatic directed cycle.
Indeed every cut must include the entire reachable set. Conversely
that set itself is a cut if neither event occurs: every edge of a
monochromatic cycle inside it would have been traversed, by the
hypothesis on $p$. In particular the uncolored test allows arbitrary directed
cycles and arbitrarily large tails.

One posting at its own owner can already change this answer. Give
$x,y,z$ respective single posting colors $a,b,b$ and tails
$\{z\},\{x,z\},\{y\}$.
Query $y$ with only $z$ posted at $b$. The search reaches all heads
but no monochromatic cycle, since it skips $y\to z$. Nevertheless
$y\star_b(z\star_a x)$ supplies the mark. The inner gate absorbs
on the $x$-row, both inner inputs are $b$ on the $y$-row, and $y$
is $b$ on the $z$-row.

The restriction on vertices belonging to cycles of several colors is
also sharp. Take the two reciprocal pairs $x\leftrightarrow_a y$ and
$x\leftrightarrow_b z$, with a separate one-post loss for each edge.
Choosing either row at $x$ leaves a monochromatic two-cycle.
Nevertheless no word omits all four losses, by
\cref{core:thm:additive-greatest-minorant}. At an uncolored query
no closed cut can avoid both pairs: every remaining head points to
$x$, or already is $x$, and whichever row is chosen at $x$ returns
to it in the same color.

When there is no monochromatic directed cycle anywhere, the cut
also tests a \emph{complete} guard family without quantifying over
positive colorings. Choose a nonempty set $H$ of heads and one row
$t^j$ at each. Put $B_r=\bigcup_{j\in H}(t^j_r\setminus H)$.
For a family with some positive bit, realizability is equivalent to
the following condition on every such choice: if the $B_r$ are pairwise disjoint and
some $B_r$ is empty, then for every $j\in H$ some loss row at $j$
is contained in $B$ componentwise. Indeed the $B_r$ give the least proper
coloring compatible with that boundary. If a head is positive there,
the cut forbids its mark. Conversely any impossible prescribed
positive gives a cut, and remains positive on restricting its
coloring to this boundary. If every bit is zero there is no term;
otherwise union of the individual witnesses realizes the family.

\begin{corollary}[Recognition by closing the endpoints of an edge]
\label{core:cor:monochromatic-boundary-recognition}
Suppose each headed input has one loss, the graph has no
monochromatic directed cycle, and some bit is positive. For every
edge $j\to u$ and every owner $k$, close $\{j,u\}$ under two rules.
Include every vertex posted at $k$ by an included head. Include a vertex
whenever included heads post it at two distinct colors. The
family is realizable exactly when every such closure reaches an
unheaded input. In particular full recognition takes polynomial time.
\end{corollary}
\begin{proof}
A proper boundary cannot enclose an edge: if a row at $j$ posts in $H$, its
head is positive at the boundary, and otherwise it is suppressed
there. For a chosen missing owner $k$, $B_k=\varnothing$ is the
first closure rule, and disjointness of the $B_r$ is the second.
Hence the least closed set containing the edge is a forbidden set
of heads precisely when it avoids the unheaded inputs. There is one
finite closure computation per edge and owner.
\end{proof}

Specialize now to heads with a designated owner $c(j)$: every loss
at $j$ posts its tail $T_j$ entirely at $c(j)$.
With at most two owners among the heads and no same-owner edge,
there is also a polynomial test on explicitly listed minimal tails,
with no bound on their number per head. For each tail $T_j$, each
headed $u\in T_j$, and each tail $T_u$, require $T_j\cap T_u\ne
\varnothing$. Any shared input is unheaded: a headed one would
receive an edge from its own owner. Thus a violating cut cannot have
an internal chosen edge $j\to u$ and satisfy the pair condition,
since its two tails would conflict outside. Conversely disjoint
tails give the cut $\{j,u\}$ with a proper boundary. Head $j$ is
positive there since its postings at $c(j)$ are just
$T_j\setminus\{u\}$, which contains no minimal tail at $j$.
Throughout, a family with all bits zero is rejected separately.

On prime lists alone, this two-owner pair test is polynomial at
width two and coNP-complete at width three. To find disjoint tails
with $u\in T_j$, choose a prime $G$ at $j$ containing $u$; all
primes here are viewed as sets of inputs on a facet seeing the
head's owner. Seek a set $X$ meeting every prime at $j$, with
$X\cap G=\{u\}$, whose complement meets every prime at $u$.
Minimizing the two hitting sets yields disjoint minimal tails and
preserves $u$ in the first. Conversely minimality of $T_j$ gives
such a $G$. Each prime gives one clause of its own width in the
membership bits of $X$; the prescribed intersection only adds unit
clauses. Thus width two gives a polynomial number of 2SAT tests.

For hardness, use just two heads $j,u$ at different owners, and
two unheaded literal inputs for every variable of a 3CNF formula.
At $j$ use the primes $\{u\}$, the literal pair of each variable,
and the sets of literals in each clause. At $u$ use just the
literal pairs. Remove redundant supersets to make the lists
canonical, and give the facets hiding the respective owners the
empty guard. The sole conditional dependency is $j\to u$. The
first hitting set and its disjoint counterpart choose opposite
literals of each variable; the first must meet every clause.
Hence an obstructing pair exists exactly for a satisfiable formula.

\begin{theorem}[Two losses per head suffice for coNP-completeness]
\label{core:thm:two-monochromatic-losses}
At each fixed $K\ge3$, full recognition is coNP-complete even with at
most two monochromatic losses per head, no edge between heads of the
same owner, and at most two exclusions per prime. Both the minimal
losses and the canonical primes may be supplied.
\end{theorem}
\begin{proof}
The boundary criterion gives a certificate of failure: list the heads,
one loss per head and a head positive at the resulting proper boundary.
With prime lists alone, a proposed minimal loss can be verified by
testing it and its one-post deletions. The family with no positive bit
is recognized directly.

Reduce 3SAT to failure of recognition. Use three owners $a,b,c$.
Give each Boolean variable two unheaded inputs for its two literals.
There are three control heads, one per owner. Call a further head
mandatory if the construction below requires every violating cut to
contain it. Each control has one tail containing the other two
controls and all mandatory heads of a different owner. The $b$-control
also posts all optional heads. Every tail at a noncontrol head contains
the two controls of owners different from its own. In addition it
posts the one input specified by the following choices.

For each variable introduce two mandatory heads, one at $a$ and one
at $b$. Each chooses either literal input of the variable.
For each clause with literals $\ell_1,\ell_2,\ell_3$, introduce a
mandatory $a$-head whose two choices are the input for $\ell_1$ and
an optional $c$-head $z$. The two choices at $z$ are optional
$a$-heads $w_2,w_3$. The sole choice at $w_j$ is the input for
$\ell_j$. Here a choice denotes one possible tail in the loss list.
All tails are monochromatic and no edge joins heads of the same
owner. On a facet seeing the loss color, a prime excludes a common
posting or excludes the private postings of both alternative tails.
A head with one tail has only single-literal primes; on the facet
hiding the color the empty guard suffices. Thus both presentations
have polynomial size.

A violating cut contains an edge of one of its chosen tails: otherwise
those tails alone suppress every selected head at the boundary.
The endpoints have distinct owners. Both post the third control, which
must therefore lie inside the cut. These three heads now force every
control inside: the two heads of owners different from a missing
control both post it. The controls in turn force every mandatory head,
since two controls of different owners post it. The opposite-color
choices at each variable pair must choose different literal inputs.
Declare the $a$-choice true. In a clause the first choice either
posts a true literal at $a$ or selects $z$. In the latter case $z$
must lie inside: outside it would also be posted at $b$ by the
control. The same argument puts the chosen $w_j$ inside, and its
$a$-posting makes $\ell_j$ true.

Conversely, given a satisfying assignment take all controls and
mandatory heads, choosing according to the assignment. In each clause
use a true literal, adjoining $z$ and its selected $w_j$ precisely
when using one of the last two literals. The boundary posts omitted
optional heads at $b$, true literal inputs at $a$, and false ones
at $b$. It is consistent and omits $c$. All controls are positive
there since their tails contain other controls. This gives the
violating cut.
\end{proof}

\begin{proposition}[A single posting missed by every boundary cut]
\label{core:prop:monochromatic-cut-obstruction}
For every $K\ge3$ a seven-input family with one monochromatic loss per
head passes all acyclic boundary cuts, even when every owner may occur
on the boundary. It has idempotent pure-owner closures and admissible
singleton profiles. Nevertheless its losses force a singleton when
just one input is posted, where the family prescribes another winner.
Only one pair of inputs forms a monochromatic cycle.
\end{proposition}
\begin{proof}
Use distinct owners $d,\alpha,\beta$. In the table a string denotes
a set of input names; upper- and lower-case names are distinct.
\[
\begin{array}{c|ccccccc}
j&A&a&X&B&Y&b&Z\\ \hline
c(j)&\alpha&\alpha&d&\beta&d&\beta&d\\
T_j&XBbZ&BbZ&ABY&AaYZ&aXB&AXZ&Ab
\end{array}
\]
The two $d$-heads $X,Y$ post one another, so their fold lies below
the family. There is no other same-owner edge. Each pure closure
adds its heads in one step: adding $X$ requires $Y$ already present,
and conversely; $Z$ is posted by neither. The $d$-profile requires
$A,a,B,b$ and at least one of $X,Y$. Each foreign profile requires
all inputs except the heads of its own owner; unused owners require
all inputs. Thus every positive rainbow has a common input in $XY$.

A consistent boundary must include every common posting of two
included heads of different owners. The following implications
list the unordered edges between different owners. Each pair on
the left forces the indicated pair by this one rule, and the rows
are ordered so that each implication leads back to $XY$.
\[
\begin{array}{c|c@{\qquad}c|c}
 XB&XY & AX&XB\\
 Ab,\ Xb&AX & AZ,\ bZ&Ab\\
 AB,\ BZ&AZ & ab,\ aZ&bZ\\
 aB&aZ & aY,\ BY&aB
\end{array}
\]
Consequently a set enclosing any edge either contains the two-cycle
$X,Y$ or has inconsistent boundary. This verifies every acyclic cut.

Write $P=\{X\}$. Omission of all $s$-rows at a set $L$ of $s$-heads
whose tails avoid $L$ forces $\Phi_s(V\setminus L)=1$: post everything
outside $L$ at $s$. Each head in $L$ is suppressed even with its
menu free, so only the singleton $s$ can result. This applies to
$Aa$ at $\alpha$, $Bb$ at $\beta$, and to $XZ,YZ,Z$ at $d$.
It rules out absorption on all four foreign rows. For example, an
$\alpha$-win on $T_A=XBbZ$ would have to meet simultaneously the
forced $d$-coalition excluding $XZ$ and the $\beta$-coalition
excluding $Bb$. The same argument covers $T_a,T_B,T_b$, using their
indicated foreign and $d$ parts.

It remains to exclude absorption on the three $d$-rows under the
assumption $\Phi_d(P)=0$. We give the smaller sets of rows needed
in each step. Whenever absorption on all retained rows is excluded,
a gate acts by union there, so both children keep the omissions.
One child keeps a same-owner winning and losing pair of tests:
at that owner's gate both lose and one wins, and at a foreign gate
both win and one loses. Likewise one child retains wins at distinct
owners. Following these choices must end at a projection.
\begin{enumerate}
\item Even with just the rows at $A,b$, no word wins the $d$-test on
$T_Z=Ab$. Its intersections with $T_A,T_b$ are respectively $b,A$,
excluded by the opposite row's forced profile. Descent would select
one of the forbidden heads $A,b$.
\item With just the rows at $a,B$, a word losing on $P$ cannot win
on $T_Y=aXB$. Its intersections with $T_a,T_B$ are $B,a$,
respectively. The opposite row again excludes absorption, and the
only projections in $T_Y\setminus P$ are $a,B$.
\item Retain the rows $A,a,B,Y$. A $d$-win on $T_X=ABY$ excludes
absorption on $A,a$, since their tails meet $T_X$ only at $B$.
The second step excludes absorption on $Y$. If $B$ were atomic,
the word would win on both $T_B$ at $\beta$ and $T_X$ at $d$.
Retain only $A,a,Y$. The two winning coalitions intersect in $AY$,
disjoint from both $\alpha$-tails, and $T_B\cap T_Y=\{a\}$, excluded
by the $\alpha$-row at $a$. Descent with the two wins would end at
$A$ or $Y$, both forbidden. Thus $B$ cannot absorb either. Descent
with the original $d$-win and loss would end at $A,B$ or $Y$.
\end{enumerate}
On all seven rows, therefore, a word losing on $P$ could be atomic
nowhere. Descending with that loss would give a bounded projection,
although every input heads a loss. Hence every bounded word gives
the $d$-singleton when just $X$ is posted at $d$. The family
prescribes the mark of $A$ there. By \cref{core:lem:singleton-profile-absorption},
at any valuation each bounded word must either return $\{d\}$ or
include the entire menu at $X$.
\end{proof}

\begin{corollary}[The activation principal from the surviving inputs]
\label{core:cor:monochromatic-principal-box}
Let $D$ be the greatest word omitting the prescribed losses, or $0$ if
there is none. Assume the two cycle hypotheses of
\cref{core:thm:monochromatic-head-cut} for a proper coloring $p$.
Let $A$ consist of the inputs $j$ whose mark is possible at
$p|_{V\setminus\{j\}}$. For every proper coloring $b$,
\[
             b\preceq_D p\quad\Longleftrightarrow\quad
                    b|_A\subseteq p|_A .
\]
If $A\ne\varnothing$, a single word omitting the losses has support
exactly $A$ and supplies all these marks. Thus the tests of
\cref{core:thm:monochromatic-head-cut} compute the whole principal.
\end{corollary}
\begin{proof}
Grow a set $J$, initially empty. Discard losses headed in $J$ and
those whose postings on $J$ are outside $p|_J$; strip $J$ from the
remaining rows. If a safe $r$-fold on a nonempty set
$T\subseteq V\setminus(J\cup p_r)$ exists for this reduced sample,
adjoin $T$ to $J$. On its own support the fold supplies every mark
at $p$ and its principal there is ordinary subcoloring:
in \cref{core:thm:one-owner-activation-order} its closure of the
empty set is empty.

Maintain a word $B$ on exactly $J$ with these two properties,
omitting the original losses. The first fold serves as $B$. For a
later fold $F$, use \cref{core:thm:block-source-extension} over $B$
with coefficient $F$ on backgrounds $q\preceq_B p|_J$ and zero
elsewhere. This extension preserves the old source bits. A loss at
a new head either lies outside that principal or is omitted by $F$
on its stripped row. All other heads outside $J\cup T$ are absent.
At $p$ each old and new source survives, giving exactly the claimed
support. The greatest coefficient at $p|_J$ contains $F$.
By \cref{core:thm:block-source-extension,core:thm:activation-faithfulness}
the new principal requires subcoloring on both $J$ and $T$;
antitonicity supplies the converse.

At termination there is no nonzero safe word on the reduced sample.
An unheaded variable would give an admissible projection; otherwise
the removal argument of \cref{core:thm:monochromatic-head-cut}
applies on the reduced graph. It gives either a safe fold on an
$r$-cycle component avoiding $p_r$, or one row at each remaining
head with no monochromatic cycle. The former would have extended
$J$, and the latter excludes every word by \cref{core:lem:ranked-posting-survivor}.
Write $R=p|_J$. By \cref{core:lem:unmarked-background}, applied after
discarding the $J$-headed losses, no word omitting the original losses
can supply any outside mark even at $R$. Thus $J=A$. If $D=0$ the
principal is constantly true. Otherwise $D$ captures all outside inputs
under $R$ by \cref{core:cor:colored-cofactor-dependence}. Every interval
predicate above $D$ positive at $p$ is positive at $R$, so
\cref{core:lem:captured-rectangle} includes the entire displayed box.
The reverse inclusion follows from $B\subseteq D$ by
\cref{core:thm:activation-faithfulness}, extending the predicates of
$B$ over the unused inputs; for $J=\varnothing$ it is automatic.
\end{proof}
